\documentclass[12pt,addpoints,answers]{exam}
\usepackage[utf8]{inputenc}
\usepackage{amsmath,amsfonts,amssymb,amsthm}
\usepackage[margin=1in]{geometry}
\usepackage{mathtools}
\usepackage{hyperref}
\usepackage{fullpage}
\usepackage{microtype}
\usepackage{xspace}
\usepackage[svgnames]{xcolor}
\usepackage[sc]{mathpazo}
\usepackage{enumitem}
\usepackage{bm}

\pagestyle{head}

% TIRET DEFINITIONS

\newlength{\saveparindent}
\setlength{\saveparindent}{\parindent}
\newlength{\saveparskip}
\setlength{\saveparskip}{\parskip}
\newcounter{ctr}
\newcounter{savectr}
\newcounter{ectr}
\newcounter{sctr}


\newenvironment{tiret}{%
\begin{list}{\hspace{2pt}\rule[0.5ex]{6pt}{1pt}\hfill}{\labelwidth=15pt%
\labelsep=5pt \leftmargin=20pt \topsep=3pt%
\setlength{\listparindent}{\saveparindent}%
\setlength{\parsep}{\saveparskip}%
\setlength{\itemsep}{0pt} }}{\end{list}}

%----------Header--------------------%
\def\course{{\sc Fundamentos y Aplicaciones de Blockchains}}
\def\term{Depto. de Computaci\'{o}n, UBA, 2do. Cuatrimestre 2025}
\def\prof{Lecturer: Juan Garay}
\newcommand{\handout}[5]{
   \renewcommand{\thepage}{\arabic{page}}
   \begin{center}
   \framebox{
      \vbox{
    \hbox to 5.78in { \hfill \large{\course} \hfill }
    \vspace{2mm}
%    \hbox to 5.78in { \hfill \large{\prof} \hfill }
%       \vspace{2mm}
       \hbox to 5.78in { {\Large \hfill \textbf{#5}  \hfill} }
       \vspace{2mm}
       \hbox to 5.78in { \term \hfill \emph{#2}}
       \hbox to 5.78in { {#3 \hfill \emph{#4}}}
      }
   }
   \end{center}
   \vspace*{4mm}
}
\newcommand{\hw}[4]{\handout{#1}{#2}{#3}{#4}{Homework #1}}

% -- For ignoring stuff -- %
\newcommand{\ignore}[1]{}

\begin{document}

%----Specs: change accordingly-----%
\newif\ifstudent % comment out false
\studenttrue 
% \studentfalse

\def\hwnum{2}
\def\issuedate{9/10/25}
\def\duedate{21/10/25, 15:00 hs} % 
\def\yourname{} % put your name here
%------------------------------%

\ifstudent
\hw{\hwnum}{\issuedate}{Student: \yourname}{Due: \duedate}%
\else
\hw{\hwnum}{\issuedate}{\prof}{Due: \duedate}%
\fi

% \ignore{

\noindent \textbf{Instructions}

\begin{itemize}
    \item Upload your solution to Campus; make sure it's only one file, and clearly write your name on the first page. Name the file \textsf{`$<$your last name$>$\_HW3.pdf.'} 
    %{\bf Important:} Make sure to tap {\bf Turn in} after you upload your solution.   
      
     If you are proficient with \LaTeX, you may also typeset your submission and submit in PDF format. To do so, uncomment the ``\%\textbackslash begin\{solution\}'' and ``\%\textbackslash end\{solution\}'' lines and write your solution between those two command lines.
    
      \item Your solutions will be graded on \emph{correctness} and
    \emph{clarity}. You should only submit work that you believe to be
    correct.
    % , and you will get significantly more partial credit if you     clearly identify the gap(s) in your solution.
    
    \item You may collaborate with others on this problem set.  However,
    you must \textbf{{write up your own solutions}} and \textbf{{list
      your collaborators and any external sources (including ChatGPT and similar generative AI chatbots)}} for each
    problem. Be ready to explain your solutions orally to a member of the course staff
    if asked.
    
    \ignore{
    \item For problems that require you to provide an algorithm, you must
    give a precise description of the algorithm, together with a proof
    of correctness and an analysis of its running time. You may use
    algorithms from class as subroutines. You may also use any facts
    that we proved in class or from the book.
    } %IGNORE
    
\end{itemize}

\noindent This homework contains \numquestions\ questions,
% \numpages\ pages
for a total of \numpoints \ points.
% and \numbonuspoints\ bonus points.

%\medskip
\newpage

%} %IGNORE

\begin{questions}


\question \textbf{Bitcoin transactions:}
         \begin{parts}
                   \part[7] Describe the mechanism used in the actual Bitcoin application that  the miners follow in order to insert transactions into a block. Would the Consistency (aka Persistence) and Liveness properties we saw in class, as well as the $V,I,R$ functions have to be modified to capture the actual mechanism? Elaborate.
                   
     \begin{solution} %Uncomment and type your solution here

        Función $V$: el objetivo es validar un bloque. En el paper no se especifican algunos detalles que puede ser útil mostrar para especificar qué es un bloque válido en la red de Bitcoin (tomo la función $V$ de la clase, la de "Public transaction ledger protocol").
        Para que un bloque sea válido en la red de bitcoin, debe cumplir las siguientes propiedades. 
        \begin{itemize}
            \item La suma total de lo que se envía en la transacción coinbase debe ser la correcta dada la recompensa actual de la red. Además, la transacción coinbase debe ser la primera del bloque.
            \item Las demás transacciones del bloque deben ser válidas. Es decir, para cada transacción:
            \begin{itemize}
                \item La firma digital es correcta respecto a quien envía la transacción.
                \item No hay double spends: la suma total de lo que se envía en esa transacción es menor o igual a los fondos de las cuentas a debitar.
                \item Si la cuenta a debitar viene de un UTXO de una coinbase transaction, entonces ese UTXO debe haber sido creado con 100 bloques de antelación. (\href{https://developer.bitcoin.org/devguide/block_chain.html#transaction-data}{referencia})
            \end{itemize}
            \item El PoW es correcto: se cumple que el hash del bloque con ese nonce es menor al target.
            \item El hash del bloque anterior es correcto.
        \end{itemize}

        Función $R$: El objetivo es leer una cadena de bloques. No hacen falta modificaciones. Al leer una cadena de bloques, solamente se verifica que sea válida.
        
        Función $I$: El objetivo es, dado un input tape, seleccionar las transacciones que se van a incluir en el nuevo. El criterio de selección es que sean válidas respecto al ledger ($\bf{x}_{\mathcal{C}}$) (y que no estén en este ledger).
        Las modificaciones que tendríamos que hacer para reflejar el protocolo de bitcoin son:
        \begin{itemize}
            \item El "input tape" se llama \textit{mempool}. Ahí se guardan las transacciones que los mineros escuchan de la red. Cada minero tiene la suya, y entre las \textit{mempools} de los mineros puede haber ligeras diferencias.
            \item El criterio de selección (además de los que se detallan arriba, que son necesarios para mantener la validez del ledger) se basa en incentivos económicos: los mineros van a elegir las transacciones que mejores comisiones tengan (\textit{transaction fees}).
            \item La "transacción neutra" que se agrega al principio de la subcadena mencionada en la definición del paper se reemplaza por la transacción coinbase, la que usan los mineros para recibir la recompensa por haber minado el bloque.
        \end{itemize}

        Propiedad de $\textit{Consistency}$: Esta propiedad muestra que, una vez que un minero (honesto) reporta una transacción $tx_0$ lo suficientemente lejos del final de la cadena (es decir, ya hubo una cantidad $k$ de bloques posteriores al bloque que tiene esta transacción $tx_0$), esa transacción va a ser reportada como estable por todos los mineros, en exactamente la misma posición. Esta es la razón por la cual las transacciones se consideran irreversibles en la red de bitcoin.
        Esta propiedad no necesita de modificaciones para reflejar el protocolo de bitcoin. Podemos agregar que en bitcoin se llaman "confirmaciones" a la cantidad de bloques que tiene esa transacción por delante. En particular, una vez que $k=6$, la transacción se considera irreversible.

        Propiedad de $\textit{Liveness}$: En el paper, esta propiedad muestra que, dada una transacción válida, si los mineros la reciben continuamente (durante $u$ rondas), eventualmente la van a reportar como estable (a más de $k$ bloques del final de la cadena).
        Para adaptarla al protocolo de bitcoin, es importante agregarle que la transacción, además de ser válida, tiene que tener un precio competitivo de $\textit{transaction fees}$. Esto se debe a que los mineros siempre van a elegir las transaccion con mejor comisión. Si la transacción está muy por debajo del promedio de transacción aceptada por los mineros, no la van a incluir en el bloque, dado que un bloque tiene un tamaño limitado. 
        Por ende, no se cumpliría liveness si no se agrega esta restricción.

    \end{solution}
    
            \part[3] Describe the purpose of a {\em coinbase} transaction. 
            
     \begin{solution} %Uncomment and type your solution here
        Las transacciones de tipo coinbase sirven para recompensar a los mineros por haber minado el bloque con éxito; quien mina el bloque tiene el derecho de agregar esta transacción a la blockchain, 
        en la que recibe esta recompensa.

        Esta transacción es distinta porque no tiene origen, y no son los mismos fondos que las comisiones que puedan cobrar los mineros. Esta es una recompensa que recibe por haber minado el bloque.
        Además, cumple el propósito de crear nuevas monedas para que circulen en la red.
        
    \end{solution}


    Nota: para escribir la solución del ejercicio 1, tomé en cuenta las diapositivas de la clase, el paper de "Bitcoin backbone" y bitcoin.org.

            \end{parts}


\newpage

\question \textbf{Proofs of work:}

\begin{parts}

\part[5] Describe the purpose and implementation of the {\bf 2x1 PoW} technique.

    \begin{solution} %Uncomment and type your solution here
    El propósito de usar la técnica 2x1 PoW es lograr que en el $\textit{1\/2 Consensus Protocol}$, que requiere que se ejecuten dos PoWs diferentes, se puedan ejecutar de a una sola vez. Esto es para prevenir que un adversario pueda ejecutar solo uno de los PoW en paralelo, pudiendo pasar su $\textit{hashing power}$ de uno al otro.
    
    Esto se debe a que si no se hiciera así, se tendrían que ejecutar dos loops distintos, sumando un total de $2q$ consultas. De este modo, se combinan ambas tareas en una sola, resultando en $q$ consultas.

    La implementación se basa en ejecutar solo un hash en vez de dos. En el $\textit{1\/2 Consensus Protocol}$ se usaban dos funciones hash, $H_0$ y $H_1$. En la función $\texttt{double\_pow}$ se usa solo una función de hash, haciendo que se tome primero un valor como el output de la función de hash "tradicional" y el otro va a ser el reverso de ese hash. 

    Se agrega una etiqueta a las tuplas existentes, para diferenciar cuál PoW es cuál.

    Notar que en la función de $\texttt{double\_pow}$ se puede lograr un PoW para uno de los dos, para los dos, o para ninguno.
        
    \end{solution}

\part[7] Let $T_1$ and $T_2$ be the target values in 2x1 PoW, and $\kappa$ the size of the hash function's output. What relation should they satisfy for the technique to work? Elaborate.

    \begin{solution} %Uncomment and type your solution here
        Sabemos que $T_1$, $T_2$ van a ser de $\kappa$ bits. Vamos a pedir que $T_1 = 2^{t_1}$ y $T_2 = 2^{t_2}$, donde $t_1$ y $t_2$ son enteros positivos.
        Para cumplir la primer prueba de double\_pow, tenemos que pedir los $\kappa - t_1$ bits más significativos del output de la función de hash deben ser $0$. Esto es porque el output de la función de hash debe ser menor a $T_1$.
        Análogo para la segunda prueba, solamente que tenemos que pedir que los bits $\textbf{menos}$ significativos sean cero. 

        Entonces, dado un hash $u$, necesitamos que sus primeros $\kappa - t_1$ bits sean cero, y sus últimos $\kappa - t_2$ bits sean cero.

        Entonces, para que estas dos condiciones sean independientes, vamos a necesitar que $\kappa - t_1 + \kappa - t_2 \le \kappa$. Esto es porque queremos que los bits que le afectan al primer PoW no afecten al segundo, 
        Simplificando queda $t_1 + t_2 \ge \kappa$. Entonces, nos quedan dos "particiones" de los bits del hash.

        Vemos que, si $T_1$ y $T_2$ no fuesen potencias de dos, terminaríamos imponiendo una condición sobre la otra partición de los bits. Por ejemplo, si $T_1 = 00101000$ ($40$), $T_2 = 00001000$ ($8$) y $u=000xxxxx$ sabemos que es válido, pues es menor que $T_1$. Ahora, $u=00100xxx$ también es válido, pero el problema es que ya estamos imponiendo una condición sobre la otra partición de los bits, los que van a corresponder con la prueba del reverso. Esta prueba va a requerir que los bits menos significativos sean cero, entonces la estamos condicionando. Entonces, si $T_1$ y $T_2$ no fuesen potencias de dos, terminaríamos imponiendo una condición sobre la otra partición de los bits. Por ende, para mantener la independencia estadística entre las pruebas, vamos a necesitar que $T_1$ y $T_2$ sean potencias de dos.


        Otra cosa a notar es que hay un "balance" entre $t_1$ y $t_2$. Si uno es muy pequeño, el otro tiene que ser muy grande, ya que la suma de ambos tiene que ser mayor a $\kappa$. Esto nos condiciona a que si queremos que un PoW sea más fácil, el otro tiene que compensar esa diferencia.


    \end{solution}
    
  \part[8] Design and argue correctness of an $\ell$x1 PoW scheme. Note that such a scheme would enable an $\ell$-parallel blockchain. What properties of the blockchain or ledger application would, if any,  benefit from a parallel blockchain? Elaborate.
  
   \begin{solution} %Uncomment and type your solution here
    
    Para diseñar un esquema $\ell$x1 PoW, vamos a necesitar particionar el hash en $\ell$ partes, de manera que todas sean estadísticamente independientes.

    Digamos que son todas partes iguales para simplificarlo. Decimos entonces que tienen tamaño $t$, haciendo que $\kappa = \ell \cdot t$.

    Luego, queremos que las $\ell$ particiones del hash sean menores a un target $T \le 2^t$. Necesitamos pedir esto, ya que sería trivial pedir que un número binario de exactamente $t$ bits sea menor a $2^t$.

    Como cada prueba mira su propia partición de bits, y ninguna partición se superpone con otra, las pruebas son estadísticamente independientes.


    Una de las propiedades que podría verse afectada es chain growth. Como los mineros trabajan en $\ell$ blockchains a la vez, cada una de ellas se vería  afectada. La red en su conjunto seguiría produciendo bloques a más o menos la misma velocidad. Pero, ahora esos bloques se van a distribuir en cada una de las $\ell$ cadenas. 

    Vemos que liveness se ve afectada. Depende de $\textit{Chain Growth}$ y $\textit{Chain Quality}$. Vemos que la cantidad de rondas $u$ que se necesitan para que todos los mineros honestos incluyan esa transacción en el bloque (con profundidad $k$) va a tener que ser $\ell$ veces mayor, puesto que se distribuye el $\textit{hashing power}$ en $\ell$ cadenas.


        
    \end{solution}
    
    \end{parts}
    
\newpage

\question \textbf{Strong consensus:} We saw how the Bitcoin backbone protocol can be used to solve the {\em consensus} problem (aka {\em Byzantine agreement}). In the {\em strong consensus} problem, the {\em Validity} condition is strengthened to require that the output value be one of the honest parties' inputs---this property is called {\em Strong Validity}. (Note that this distinction is relevant only in the case of non-binary inputs.)

\begin{parts}

\part[5] What should be the assumption on the adversarial computational power (similar to the ``Honest Majority Assumption'') for Strong Validity to hold?

    \begin{solution} %Uncomment and type your solution here
        Si tomamos $m$ como la cantidad de inputs posibles, 
        tenemos que restringir el $\textit{hashing power}$ de los adversarios por un factor de $m - 1$. Entonces, el $\textit{assumption}$ que tenemos que hacer es $t \le \frac{1-\delta}{m-1} \cdot (n-t)$. Usamos $m-1$ ya que en el caso de $m = 2$ querríamos dejarlo igual que antes.
        
        Esta cota es suficiente. Pongamos el caso de que los jugadores honestos estan divididos en los primeros $m-1$ valores, y el adversario elige el $m$-ésimo valor. Querríamos que el adversario no pueda ganarle en votos a incluso la división con menos jugadores de los honestos. Digamos que los honestos se dividen equitativamente (el peor caso, ya que así no tienen una mayoría abrumadora de votos), luego habría $\frac{n-t}{m-1}$ votos para esa división. El adversario tendría que tener al menos $\frac{n-t}{m-1}$ votos para ganarle, lo cual es imposible si $t \le \frac{1-\delta}{m-1} \cdot (n-t)$.
    \end{solution}

\part[5]  State and prove the Strong Validity lemma.


    \begin{solution} %Uncomment and type your solution here
        Dado un protocolo BA que opera sobre un conjunto de valores de entrada $V$, con $|V| = m$, y bajo la suposición de que $t \le \frac{1-\delta}{m-1} \cdot (n-t)$, asumiendo $k = O(\kappa)$ y $L = \Omega(\kappa)$, se satisface Strong Validity con probabilidad $1 - e^{-\Omega(\kappa)}$.

        Strong Validity nos dice que tenemos que el output (valor elegido) sea uno de los inputs honestos. Para asegurarnos esto usamos la suposición anterior, que nos asegura que mantenemos el hashing power del protocolo anterior.

        Como vemos en la demostración anterior de Validity, vale que $C^{\lceil k}$ tiene al menos k bloques, y luego vale Chain quality con la misma fracción de bloques adversarios.
        
    \end{solution}

\end{parts}

\newpage

\question \textbf{Smart contract programming:}  In this assignment you will create your {\bf own custom token}. Your contract should implement the public API described below:

\begin{itemize}
\item \underline{\textbf{owner:}} a public payable address that defines the contract’s ``owner,'' that is, the user that deploys the contract
\item \textbf{\textit{Transfer(address indexed from, address indexed to, uint256 value):}} an event that contains two addresses and a uint256
\item \textbf{\textit{Mint(address indexed to, uint256 value):}} an event that contains an address and a uint256
\item \textbf{\textit{Sell(address indexed from, uint256 value):}} an event that contains an address and a uint256

\item \textbf{totalSupply():} a view function that returns a uint256 of the total amount of minted tokens
\item \textbf{balanceOf(address\_account):} a view function returns a uint256 of the amount of tokens
an address owns
\item \textbf{getName():} a view function that returns a string with the token’s name
\item \textbf{getSymbol():} a view function that returns a string with the token’s symbol
\item \textbf{getPrice():} a view function that returns a uint128 with the token’s price (at which users can
redeem their tokens)

\item \textbf{transfer(address to, uint256 value):} a function that transfers \emph{value} amount of tokens
between the caller’s address and the address to; if the transfer completes successfully, the
function emits an event \emph{Transfer} with the sender’s and receiver’s addresses and the amount
of transferred tokens and returns a boolean value (\emph{true})
\item \textbf{mint(address to, uint256 value):} a function that enables \emph{only the owner} to create value new tokens and give them to address to; if the operation completes successfully, the
function emits an event \emph{Mint} with the receiver’s address and the amount of minted tokens
and returns a boolean value (\emph{true})
\item \textbf{sell(uint256 value):} a function that enables a user to sell tokens for wei at a price of \emph{600 wei per token}; if the operation completes successfully, the sold tokens are removed from
the circulating supply, and the function emits an event \emph{Sell} with the seller’s address and the
amount of sold tokens and returns a boolean value (\emph{true})
\item \textbf{close():} a function that enables \emph{only the owner} to destroy the contract; the contract’s
balance in wei, at the moment of destruction, should be transferred to the owner's address
\item \textbf{fallback} functions that enable anyone to send Ether to the contract’s account
\item \textbf{constructor} function that initializes the contract as needed

\end{itemize}
      
You should implement the smart contract and deploy it on the course's Sepolia Testnet. Your contract should be as secure and gas efficient as possible. After deploying your contract, you should buy, transfer, and sell a token in the contract. 

Your contract should implement the above API \textbf{exactly as specified}. \emph{Do not} omit implementing one of the above variables/functions/events, do not change their name or parameters, and do not add other public variables/functions. You can define other private/internal functions/variables, if necessary.

You should provide:

\begin{parts}

\part[5] A detailed description of your high-level design decisions, including (but not limited to):
\begin{tiret}
\item  What internal variables did you use?
\item What is the process of buying/selling tokens and changing the price?
\item How can users access their token balance?
\end{tiret}

    \begin{solution} %Uncomment and type your solution here
        Variables internas: $\texttt{owner}$, $\texttt{tokenBalance}$, $\texttt{\_totalSupply}$, $\texttt{price}$ y finalmente $\texttt{etherBalance}$. 
        Si el precio no necesitaba ser modificado, podría haberlo dejado como una constante, pudiendo obtenerlo mediante la función $\texttt{getPrice}$. Sin embargo, como la especificación mencionaba "cómo cambiar el precio", agregué la función $\texttt{changePrice}$, que solo puede ser llamada por el propietario del contrato.
        El mapping $\texttt{tokenBalance}$ guarda para cada dirección el saldo de la cuenta.
        El mapping $\texttt{etherBalance}$ guarda para cada dirección el monto de ether recibido por la misma. Esto permite a los usuarios enviar ether al contrato, y luego usarlo para comprar tokens. Al ser en dos pasos, los usuarios podrian tener un balance de ether y hacer varias transacciones para comprar tokens (por ejemplo, si quisieran especular con el precio, y hacer compras y ventas en diferentes momentos para ganar con la diferencia de precio).
        Luego, para retirar su ether, usan la función $\texttt{withdraw}$.

        Para comprar tokens, los usuarios usan la función $\texttt{buy}$. Esta función no estaba especificada en la API, pero la agregué por simplicidad. Esto permite a los usuarios comprar tokens con el ether que ya habían enviado al contrato, que se guarda en $\texttt{etherBalance}$. Luego, el usuario le compra los tokens al contrato directamente. Esto es una decisión de diseño. Consideré que era más simple que le compren al contrato (asumiendo que el contrato tiene tokens en su balance, generados por el propietario previamente). 

        Para vender tokens, los usuarios llaman a $\texttt{sell}$. Esto genera un evento $\texttt{Sell}$ con la dirección del vendedor y la cantidad de tokens vendidos. Luego, los usuarios pueden retirar su ether mediante la función $\texttt{withdraw}$. El hecho de que hayamos definido que los usuarios pueden guardar su ether en nuestro contrato permite que el withdraw no se haga en el momento de la llamada a $\texttt{sell}$, sino que se pueda hacer después. Esto nos da más seguridad en caso de que se revierta la ejecución.

        Los usuarios acceden a su balance de tokens mediante la función $\texttt{balanceOf}$. Además, acceden a su balance de ether mediante la función $\texttt{etherBalanceOf}$.

    \end{solution}

\part[5] A detailed gas evaluation of your implementation, including:
\begin{tiret}
\item The cost of deploying and interacting with your contract.
\item Techniques to make your contract more cost effective.
\end{tiret}

    \begin{solution} %Uncomment and type your solution here
        Gas de las interacciones:
        \begin{itemize}
            \item Deploy: 
            \item buy:
            \item sell:
            \item withdraw:
        \end{itemize}


        
        Para hacer que el contrato gaste menos gas, se podrían hacer ciertos recortes. Por ejemplo, en vez de permitir que se guarde el balance de ether, en otro arreglo, podríamos hacer que los usuarios envíen el ether directamente en la llamada a $\texttt{buy}$, y que al vender se les trasfiera el ether resultante directamente a su dirección. El costo de gas sería menor ya que tendría menos escrituras y lecturas de memoria, aunque el contrato se vuelve menos seguro.

        
    \end{solution}

\part[5] A thorough listing of potential hazards and vulnerabilities that can occur in the smart
contract and a detailed analysis of the security mechanisms that can mitigate these
hazards.

    \begin{solution} %Uncomment and type your solution here
        Los principales problemas de seguridad están relacionados con los privilegios que tiene el propietario del contrato. Para empezar, el dueño del contrato puede retirar todo el ether del contrato, lo que haría que los usuarios no puedan vender sus tokens. Además, el dueño puede manipular arbitrariamente el precio de los tokens. 
        
        Por estas razones el contrato no es del todo justo para los usuarios, que corren el riesgo de perder su ether si el contrato se cierra.

        Para mitigar estos problemas, se podría hacer que el dueño del contrato no pueda retirar el ether de otros usuarios, solo de su propia cuenta.
        También que no pueda cambiar el precio de los tokens arbitrariamente. Hay varias soluciones a esto, una es que el dueño del contrato pueda cambiar el precio, pero que haya un delay entre que lo propone y que se hace efectivo. Si no, podria pasar que el dueño del contrato vea que alguien esta por vender tokens, y le baje el precio para que, al vender, el usuario reciba menos ether. También se podría hacer que los usuarios que quieren comprar/vender sus tokens pongan un precio mínimo/máximo que aceptan (en esa transacción).

    \end{solution}

\part[5] The transaction history of the deployment of and interaction with your contract.

    %\begin{solution} %Uncomment and type your solution here
        
    %\end{solution}

\part[5] The code of your contract.

    %\begin{solution} %Uncomment and type your solution here
        
    %\end{solution}

\end{parts}

\newpage


~\\

\end{questions}

\end{document}
